% Options for packages loaded elsewhere
\PassOptionsToPackage{unicode}{hyperref}
\PassOptionsToPackage{hyphens}{url}
\PassOptionsToPackage{dvipsnames,svgnames,x11names}{xcolor}
\PassOptionsToPackage{space}{xeCJK}
\documentclass[
  a4paper,
  10pt]{article}
\usepackage{xcolor}
\usepackage{amsmath,amssymb}
\setcounter{secnumdepth}{5}
\usepackage{iftex}
\ifPDFTeX
  \usepackage[T1]{fontenc}
  \usepackage[utf8]{inputenc}
  \usepackage{textcomp} % provide euro and other symbols
\else % if luatex or xetex
  \usepackage{unicode-math} % this also loads fontspec
  \defaultfontfeatures{Scale=MatchLowercase}
  \defaultfontfeatures[\rmfamily]{Ligatures=TeX,Scale=1}
\fi
\usepackage{lmodern}
\ifPDFTeX\else
  % xetex/luatex font selection
  \setmainfont[]{Times New Roman}
  \setmonofont[]{Menlo}
  \setmathfont[]{STIX Two Math}
  \ifXeTeX
    \usepackage{xeCJK}
    \setCJKmainfont[]{Songti SC}
  \fi
  \ifLuaTeX
    \usepackage[]{luatexja-fontspec}
    \setmainjfont[]{Songti SC}
  \fi
\fi
% Use upquote if available, for straight quotes in verbatim environments
\IfFileExists{upquote.sty}{\usepackage{upquote}}{}
\IfFileExists{microtype.sty}{% use microtype if available
  \usepackage[]{microtype}
  \UseMicrotypeSet[protrusion]{basicmath} % disable protrusion for tt fonts
}{}
\setlength{\emergencystretch}{3em} % prevent overfull lines
\providecommand{\tightlist}{%
  \setlength{\itemsep}{0pt}\setlength{\parskip}{0pt}}
\usepackage[a4paper,top=24mm,bottom=25mm,left=23mm,right=23mm,headheight=15pt]{geometry}
\usepackage{amsmath,amssymb,mathtools,bm}
\usepackage{booktabs,longtable,array}
\usepackage{graphicx}
\usepackage{float}
\usepackage{caption}
\usepackage{enumitem}
\usepackage{fancyhdr}
\usepackage{titlesec}
\usepackage{mdframed}
\usepackage{xcolor}
\usepackage{xurl}
\usepackage{hyperref}
\graphicspath{{../../}}

\definecolor{PaperBlue}{HTML}{000000}
\definecolor{PaperRule}{HTML}{B7C6D5}
\definecolor{PaperShade}{HTML}{F4F7FA}
\definecolor{PaperText}{HTML}{000000}

\hypersetup{
  colorlinks=true,
  linkcolor=PaperBlue,
  urlcolor=PaperBlue,
  citecolor=PaperBlue,
  pdfborder={0 0 0}
}

\color{PaperText}
\linespread{1.08}
\setlength{\parindent}{1.5em}
\setlength{\parskip}{0.25em}
\setlength{\emergencystretch}{3em}
\setlist{itemsep=0.2em,topsep=0.45em}
\setcounter{secnumdepth}{3}
\setcounter{tocdepth}{3}

\titleformat{\section}
  {\Large\bfseries\color{PaperBlue}}
  {\thesection}{0.7em}{}
\titleformat{\subsection}
  {\large\bfseries\color{PaperText}}
  {\thesubsection}{0.65em}{}
\titleformat{\subsubsection}
  {\normalsize\bfseries\color{PaperText}}
  {\thesubsubsection}{0.6em}{}
\titlespacing*{\section}{0pt}{2.2ex plus 0.5ex}{0.9ex}
\titlespacing*{\subsection}{0pt}{1.8ex plus 0.4ex}{0.7ex}
\titlespacing*{\subsubsection}{0pt}{1.5ex plus 0.3ex}{0.5ex}

\captionsetup{
  font=small,
  labelfont=bf,
  textfont=it,
  justification=centering,
  skip=6pt
}

\pagestyle{fancy}
\fancyhf{}
\fancyhead[L]{\small\nouppercase{\leftmark}}
\fancyhead[R]{\small Selene's Blog}
\fancyfoot[C]{\small\thepage}
\renewcommand{\headrulewidth}{0.4pt}
\renewcommand{\footrulewidth}{0pt}
\fancypagestyle{plain}{
  \fancyhf{}
  \fancyhead[R]{\small Selene's Blog}
  \fancyfoot[C]{\small\thepage}
  \renewcommand{\headrulewidth}{0.4pt}
}

\renewenvironment{quote}
  {\begin{mdframed}[
      linewidth=0.7pt,
      linecolor=PaperRule,
      backgroundcolor=PaperShade,
      leftline=true,
      rightline=false,
      topline=false,
      bottomline=false,
      innerleftmargin=10pt,
      innerrightmargin=8pt,
      innertopmargin=7pt,
      innerbottommargin=7pt,
      skipabove=8pt,
      skipbelow=10pt
    ]\itshape}
  {\end{mdframed}}

\newenvironment{theorembox}[1]
  {\par\medskip\noindent\textbf{Theorem (#1).}\quad}
  {\par\medskip}

\newenvironment{lemmabox}[1]
  {\par\medskip\noindent\textbf{Lemma #1.}\quad}
  {\par\medskip}

\newenvironment{proofbox}
  {\par\medskip\noindent\textit{Proof.}\quad}
  {\hfill$\square$\par\medskip}

\AtBeginDocument{\allowdisplaybreaks[2]}
\usepackage{bookmark}
\IfFileExists{xurl.sty}{\usepackage{xurl}}{} % add URL line breaks if available
\urlstyle{same}
% fallback for those not using the hyperref driver hyperxmp:
\makeatletter
\@ifundefined{xmpquote}{\newcommand{\xmpquote}[1]{#1}}{}
\makeatother
\hypersetup{
  pdftitle={Two Proofs of the Matrix-Tree Theorem},
  pdfauthor={Selene},
  colorlinks=true,
  linkcolor={black},
  filecolor={Maroon},
  citecolor={Blue},
  urlcolor={black},
  pdfcreator={LaTeX via pandoc}}

\title{Two Proofs of the Matrix-Tree Theorem}
\author{Selene}
\date{June 25, 2026}

\begin{document}
\maketitle

This chapter introduces the Matrix-Tree Theorem and presents two proofs,
using induction and the Cauchy-Binet formula, to show how determinants
can be used to count spanning trees in a graph.

We denote by \(A[i]\) the matrix obtained from \(A\) by deleting the
\(i\)-th row and the \(i\)-th column.

\section{Kirchhoff's Matrix-Tree
Theorem}\label{kirchhoffs-matrix-tree-theorem}

\begin{theorembox}{Kirchhoff's Matrix-Tree Theorem}

The number of spanning trees of a graph \(G\) is given by

\[
\tau(G)=\det(L_G[i]),
\]

where \(i\) can be any vertex.

\end{theorembox}

Here \(L_G\) denotes the Laplacian matrix of \(G\).

\section{Fact 1}\label{fact-1}

Let \(A \in \mathbb{R}^{n \times n}\), and let \(E_{ii}\) be the matrix
whose \((i,i)\) entry is \(1\) and all other entries are \(0\). Then

\[
\det(A + E_{ii})
=
\det(A) + \det(A[i]).
\]

Consider the permutation expansion definition of the determinant:

\[
\det(A=(a_{ij}))
=
\sum_{\pi}
\operatorname{sgn}(\pi)
a_{1\pi(1)}
a_{2\pi(2)}
\cdots
a_{n\pi(n)}.
\]

This fact is not difficult to see. We replace the \((i,i)\) entry by
\(a_{ii}+1\), so we obtain the original determinant plus one extra term.

This extra term corresponds exactly to the sum over all permutations on
the set \([n]\setminus\{i\}\), namely the determinant obtained by
deleting the \(i\)-th row and the \(i\)-th column. Hence this term is
precisely \(\det(A[i])\).

\section{First Proof}\label{first-proof}

Now we begin the first proof of Kirchhoff's Matrix-Tree Theorem.

Our first proof proceeds by induction on both the number of vertices and
the number of edges of the graph \(G\).

If \(G\) is an empty graph with two vertices, then

\[
L_G=
\begin{bmatrix}
0 & 0\\
0 & 0
\end{bmatrix}.
\]

Hence \(L_G[i]=[0]\), and

\[
\det(L_G[i])=0,
\]

which agrees with the theorem.

Throughout the proof:

\begin{itemize}
\tightlist
\item
  \(\tau(G)\) denotes the number of spanning trees of \(G\);
\item
  \(G-e\) denotes the graph obtained by deleting the edge \(e\);
\item
  \(G/e\) denotes the graph obtained by contracting the edge \(e\), that
  is, merging its two endpoints into a single vertex.
\end{itemize}

If \(i\) is an isolated vertex, then \(G\) has no spanning trees, and
the \(i\)-th row and the \(i\)-th column of \(L_G\) are entirely zero.
Therefore

\[
\det(L_G[i])=\det(L_{G-i})=0.
\]

Otherwise, we may assume there exists an edge \(e=(i,j)\) adjacent to
\(i\). For any spanning tree \(T\), either \(e \in T\) or
\(e \notin T\), so

\[
\tau(G)=\tau(G/e)+\tau(G-e).
\]

The first term reduces the number of vertices by \(1\), while the second
term reduces the number of edges by \(1\), so this gives us an induction
framework.

First, we establish the relation between \(L_G\) and \(L_{G-e}\).
Observe that

\[
L_G[i]=L_{G-e}[i]+E_{jj}.
\]

That is, after deleting the edge \(e\), the only difference between the
resulting Laplacian matrix and \(L_G[i]\) lies in the degree of vertex
\(j\). By Fact 1, we obtain

\[
\begin{aligned}
\det(L_G[i])
&= \det(L_{G-e}[i]+E_{jj}) \\
&= \det(L_{G-e}[i])+\det(L_{G-e}[i,j]) \\
&= \det(L_{G-e}[i])+\det(L_G[i,j]).
\end{aligned}
\]

Here \(L_G[i,j]\) denotes the matrix obtained by deleting both the
\(i\)-th and \(j\)-th rows and columns. The last equality holds because
once the \(i,j\) rows and columns are removed, there is no remaining
difference between \(L_G\) and \(L_{G-e}\).

Next, we establish the relation between \(L_G\) and \(L_{G/e}\). Suppose
we contract vertex \(i\) into vertex \(j\). Then

\[
L_{G/e}[j]=L_G[i,j].
\]

Hence

\[
\begin{aligned}
\det(L_G[i])
&= \det(L_{G-e}[i])+\det(L_{G/e}[j]) \\
&= \tau(G-e)+\tau(G/e) \\
&= \tau(G).
\end{aligned}
\]

The second equality follows from the induction hypothesis, completing
the first proof.

\section{Fact 2: Cauchy-Binet
Formula}\label{fact-2-cauchy-binet-formula}

Let

\[
A \in \mathbb{R}^{n \times m},
\qquad
B \in \mathbb{R}^{m \times n},
\]

where \(m \ge n\). Define:

\begin{itemize}
\tightlist
\item
  \(A_S\): the submatrix of \(A\) consisting of the columns indexed by
  \(S\subseteq[m]\);
\item
  \(B_S\): the submatrix of \(B\) consisting of the rows indexed by
  \(S\subseteq[m]\).
\end{itemize}

Let \(\binom{[m]}{n}\) denote the collection of all subsets of \([m]\)
of size \(n\). Then

\[
\det(AB)
=
\sum_{S \in \binom{[m]}{n}}
\det(A_S)\det(B_S).
\]

\section{Second Proof}\label{second-proof}

Since

\[
L_G = \sum_{(i,j) \in E}(e_i - e_j)(e_i-e_j)^T,
\]

we may write

\[
L_G=BB^T,
\]

where \(B \in \mathbb{R}^{n \times m}\) is an incidence matrix of \(G\).
Each edge \((i,j)\) corresponds to a column vector

\[
e_i-e_j.
\]

Since \(L_G=BB^T\), this once again proves that \(L_G\) is positive
semidefinite.

If \(B[i]\) denotes the matrix obtained by deleting the \(i\)-th row of
\(B\), then

\[
L_G[i]=B[i]B[i]^T.
\]

We further let \(B_S[i]\) denote the matrix obtained from \(B[i]\) by
keeping only the columns indexed by \(S\subseteq E\).

We also need the following lemma.

\begin{lemmabox}{2}

For \(S \subseteq E\) with \(\lvert S\rvert = n-1\),

\[
\left|\det(B_S[i])\right|
=
\begin{cases}
1, & \text{if } S \text{ is a spanning tree},\\
0, & \text{otherwise}.
\end{cases}
\]

\end{lemmabox}

With this lemma, the second proof of the Matrix-Tree Theorem becomes
extremely short:

\[
\begin{aligned}
\det(L_G[i])
&= \det(B[i]B[i]^T) \\
&= \sum_{S\in\binom{E}{n-1}}
   \det(B_S[i])^2 \\
&= \tau(G).
\end{aligned}
\]

Here:

\begin{itemize}
\tightlist
\item
  the second equality follows from the Cauchy-Binet formula;
\item
  the third equality follows from Lemma 2.
\end{itemize}

\section{Proof of Lemma 2}\label{proof-of-lemma-2}

\begin{proofbox}

Assume without loss of generality that the edges in \(B_S[i]\) are
oriented arbitrarily. That is, we may replace the column corresponding
to an edge \((i,j)\) from \(e_i-e_j\) to \(e_j-e_i\).

This only changes the sign of the determinant, so it does not change the
absolute value.

If \(S\subseteq E\), \(\lvert S\rvert=n-1\), and \(S\) is not a spanning
tree, then \(S\) necessarily contains a cycle. We may orient the edges
along this cycle.

If we sum the columns corresponding to this cycle, we obtain the zero
vector. Hence the columns of \(B_S[i]\) are linearly dependent, and
therefore

\[
\det(B_S[i])=0.
\]

Now assume that \(S\) is a spanning tree. We prove

\[
\left|\det(B_S[i])\right|=1
\]

by induction on \(n\).

\subsection{Base Case}\label{base-case}

If \(n=2\), then

\[
B_S=
\begin{bmatrix}
1\\
-1
\end{bmatrix}.
\]

After deleting one row, we get

\[
B_S[i]=[\pm 1],
\]

and therefore

\[
\left|\det(B_S[i])\right|=1.
\]

\subsection{Inductive Case}\label{inductive-case}

Assume the lemma holds for trees with \(n-1\) vertices.

Let \(j\ne i\) be a leaf vertex of the tree, and let \((k,j)\) be the
unique edge adjacent to \(j\). Such a leaf vertex exists because every
tree with at least two vertices has at least two leaves.

Permute rows and columns so that:

\begin{itemize}
\tightlist
\item
  the edge \((k,j)\) becomes the last column;
\item
  the vertex \(j\) becomes the last row.
\end{itemize}

These permutations may change the sign of the determinant, but they do
not affect its absolute value.

Since \(j\) is a leaf, the last row has only one nonzero entry.
Expanding the determinant along the last row gives

\[
\left|\det(B_S[i])\right|
=
\left|
\det(B_{S-\{(k,j)\}}[i,j])
\right|.
\]

By the induction hypothesis,

\[
\left|
\det(B_{S-\{(k,j)\}}[i,j])
\right|
=1.
\]

Therefore,

\[
\left|\det(B_S[i])\right|=1.
\]

This completes the proof of Lemma 2, and hence the second proof of
Kirchhoff's Matrix-Tree Theorem.

\end{proofbox}

\end{document}
